\section{Discussion}\label{sec:discussion}

\textbf{Where MOEA/D stands.} The answer to RQ1 depends on the problem and on how the baselines are treated. With fixed budgets and no tuning (Sections 6.2 and 6.7) the MOEA/D variants are four to six ranks ahead of NSGA-II, GRASP and VNS everywhere; with a tuning of the same size for every method and the same wall-clock time (Section 6.12) the advantage survives on the GDP and, at equal time, on the CDP, disappears with five times the time on the CDP, and reverses on the unconstrained MDP, where a tabu search or a VNS with the front as incumbent and shake-and-local-search cycles is better. Decomposition with a repair step and a pool of operators pays where there are side constraints to satisfy while exploring; where there are none, a good single-objective search applied to both ends of the front is hard to beat --- the same reason that made the MaxMin anchor of Section 6.6 useful. This agrees with Colmenar et al., who found a tabu search to be the best method for the MDP. The hypervolume-based SMS-EMOA, added as a further baseline, behaves like NSGA-II: tuned, it ends 0.11--0.37 below the tuned MOEA/D in every cell (Section 6.12), so the conclusion does not hinge on the choice of the Pareto-based MOEA.

\textbf{Where the performance comes from.} Taken together, Sections 6.3--6.5 say that on small and medium instances the benefit of the MOEA/D framework is the \emph{pool}: eight complementary operators mixed at random inside a decomposition that scales its objectives. Choosing among the operators online adds nothing measurable there, and the most sophisticated selectors (scalar-reward bandits) subtract. We see two reasons. The operators act on different aspects of a solution (global restart, bottleneck repair, recombination, perturbation), and their value depends on the state of the subproblem and on the other operators; uniform mixing keeps all of them alive, while any reward-greedy policy concentrates on whichever operator looks best now. The ablations make ``complementary'' precise: no operator suffices alone, MaxMinSwap cannot be removed (on the GDP neither can SubsetCrossover and RandomSwap2$\times$2), the other operators are individually dispensable, and the uniform pool is not damaged by three deliberately useless ones. The last point is a property of elitist replacement: a child replaces a neighbour only if it is not worse, so a useless operator wastes effort but cannot harm the population, and at this budget the wasted effort costs less than the noise. The second reason is that the reward signals available in a bi-objective archive are weak or myopic: hypervolume gain saturates early, scalarised improvement favours local moves. Section 6.11 qualifies this: where the search is far from convergence (2,000 elements) or the constraint is tight (GDP cost factor 0.2), the hypervolume reward does not saturate in time and concentrating on BiObjBestSwap and MaxMinSwap pays, by up to +0.15 normalised hypervolume. The tight cost budget is the case in which a third of the children need repair and the search is farthest from convergence at the standard budget (Section 6.11), which is why reallocating effort to the two productive operators pays there; we could show this for groups of instances, not for individual instances. So the pool and the selector are not alternatives: the pool provides the diversity, and a selector earns its keep when the budget is short relative to what the search has to do. This does not show that adaptive selection cannot help in other pools or with other policies; it shows where the standard bandit recipes do and do not.

\textbf{Implications for practice.} (i) Scale the objectives inside the scalarisation before anything else; the effect grows with n and can reverse conclusions about baselines. (ii) Use a diverse pool with uniform selection as the default for small and medium instances; for instances with thousands of elements or with tight constraints, and for short budgets, test an HV-reward bandit against that default under replication --- it may be worth tens of percent of hypervolume. (iii) Check the ends of the front explicitly; HV-type indicators can hide a large MaxMin gap, and a cheap single-objective intensification narrows it. (iv) Tune the baselines with the same protocol as the proposed method before claiming an advantage: it changed the answer on the MDP and removed the advantage at five times the time on the CDP. (v) For the unconstrained MDP use a tabu search or VNS with the front as incumbent; keep the MOEA/D framework for the capacity- and budget-constrained variants. (vi) State the size regime of the constrained bi-objective problems: without a cardinality bound the MaxSum objective rewards growth. (vii) Compare objective values with a relative tolerance in any success-based reward or counter: the same set of elements evaluated in another order differs by about 10$^{-}$$^{1}$$^{5}$, which silently turned duplicates into successes in our first runs and distorted the selectors that count successes.

\textbf{Limitations.} Budgets of the main experiment are fixed by generations, not by time, and its baselines are untuned and have small budgets; Section 6.12 repeats the comparison with tuned baselines and with the tabu search and VNS of \citet{colmenar2018}, who tuned with up to 500 experiments per algorithm (our tuning: about 56 runs per method and cell, eight training instances), and our tabu search and VNS are reimplementations from their description, not their code (none exists for CDP/GDP to our knowledge); the fronts of the CDP and GDP refer to the size regime that the capacity constraint induces (Section 6.13), so $\mathrm{HV}_{\mathrm{norm}}$ is not a fraction of the true Pareto front; replicate counts are three (five for small instances) and the largest instances (MDG-b n = 2000, MDG-c n = 3000) are covered by a five-instance sample only; the budget-relaxation analysis covers GDP only and n = 500 with twenty instances, and the time-matched comparison covers every instance with n $\leq$ 150 but only 20 CDP and 40 GDP instances with n = 500 (the tuned-baseline comparison uses 20 and 32, and the MDP instances with 100 $\leq$ n $\leq$ 150 and with n = 500 of the GKD-c and MDG-b families only); the study uses one hardware platform; no indicator-based or other decomposition-based MOEA besides SMS-EMOA was tested; the mechanism behind the effect of the tight cost budget is established for groups of instances only; the raw-objective catalogue of Section 6.1 was not recomputed after the correction of Section 4.3 and its two Thompson variants carry the artefact; the instance families are the standard MDPLIB/CDP/GDP benchmarks and may not represent all dispersion applications; the reference front depends on the runs included (it is frozen at submission). Selection policies richer than the ones we tested (for instance, contexts learned from the solution rather than fixed by the weight vector and the clock) remain open.
